\documentclass[11pt, a4paper]{article}
\usepackage[a4paper, top=2.5cm, bottom=2.5cm, left=2cm, right=2cm]{geometry}
\usepackage{amsmath,amssymb,amsthm,microtype}
\usepackage{hyperref}

\theoremstyle{definition}
\newtheorem{problem}{Problem}

\begin{document}

\title{\vspace{-1.5cm}\textbf{Regional Quantitative Problem-Solving Circle} \\ \Large Problem Set 1: Extremal Principles \& Monovariants (Weeks 1--2)}
\author{\textbf{Lead Architect:} Mohamed Jad Srifi}
\date{\textbf{Term:} Fall 2025 -- Spring 2026}
\maketitle

\noindent\textbf{Instructions:} Do not rely on heuristic guessing or brute-force algebra. Every problem in this set requires the identification of a structural extreme (a maximum or minimum element) or a monovariant state function. Formalize your logic using the Well-Ordering Principle or finite set properties before presenting your whiteboard defense.

\vspace{0.5cm}

\begin{problem}[Combinatorial Geometry / Extremal Distances]
Let $S$ be a configuration of $n \ge 3$ points in the Euclidean plane such that all pairwise distances between the points are mutually distinct. Each point connects a straight line segment to its strictly nearest neighbor in $S$. 
\begin{enumerate}
    \item Prove that there exists at least one pair of points that mutually connect to each other.
    \item Prove that no two line segments in this constructed network intersect.
\end{enumerate}
\end{problem}

\vspace{0.5cm}

\begin{problem}[The Sylvester-Gallai Theorem]
Let $S$ be a finite set of $n \ge 3$ points in $\mathbb{R}^2$. Assume that the points of $S$ are not all collinear (i.e., they do not all lie on a single straight line). An ``ordinary line'' is defined as a line that passes through exactly two points of $S$. 

Prove that there exists at least one ordinary line. 
\textit{(Hint: Proceed by contradiction. Consider the set of all lines connecting at least two points, and use an extremal metric regarding the distances from points to these lines.)}
\end{problem}

\vspace{0.5cm}

\begin{problem}[Infinite Lattice Harmonicity / Monovariant]
Consider an infinite two-dimensional square grid where every cell contains a strictly positive integer. The grid possesses a ``harmonic'' property: the integer in every single cell is exactly equal to the arithmetic mean of the integers in its four orthogonally adjacent neighboring cells (above, below, left, and right).

Using the Well-Ordering Principle, rigorously prove that all the integers across the entire infinite grid must be identical.
\end{problem}

\end{document}
